% ============================================================================
% METHODOLOGY — PRISM AAAI 2026 (Final Rewrite)
% Ground-truth values verified against prism/ source and config/default.yaml
% ============================================================================

\section{Methodology}
\label{sec:method}

\subsection{Problem Definition and Core Formalisms}

A Constraint Satisfaction Problem is a triple $\mathcal{P} = (\mathbf{X}, \mathbf{D}, \mathbf{C})$ where $\mathbf{X} = \{x_1, \ldots, x_n\}$ is a set of variables, $\mathbf{D} = \{D_1, \ldots, D_n\}$ specifies finite domains ($x_i \in D_i$), and $\mathbf{C} = \{c_1, \ldots, c_m\}$ is a set of constraints. A solution is a complete assignment $\theta: \mathbf{X} \to \bigcup_i D_i$ with $\theta(x_i) \in D_i$ and $c_j(\theta)$ holding for all $j$.

\textbf{Solving trajectory.} For puzzle $\mathcal{P}$, a solving trajectory is a sequence
\[
\mathcal{T} = \langle (S_0, a_1, S_1),\; (S_1, a_2, S_2),\; \ldots,\; (S_{T-1}, a_T, S_T) \rangle
\]
where $S_t = (C_t, \delta_t)$ is the solver state at step $t$ (formalized constraints $C_t$ and domain assignments $\delta_t: \mathbf{X} \to 2^{\bigcup_i D_i}$), and $a_t$ is an LLM-generated inference: a natural-language statement paired with its Z3 formalization. $\mathcal{T}$ is \textit{successful} if Z3 returns SAT at $S_T$ and the assignment matches ground truth.

\textbf{Solving paradigm.} A paradigm is a formally-structured five-tuple:
\[
P = (\mathcal{Q},\; \mathcal{O},\; \phi_{\text{pre}},\; \phi_{\text{post}},\; \mathcal{S})
\]
where $\mathcal{Q}$ specifies trigger conditions (constraint types and solver-state features that activate the paradigm); $\mathcal{O}$ is the inference operation, given as a parameterized template (e.g., $\texttt{pos}(X) = k + d$); $\phi_{\text{pre}}$ and $\phi_{\text{post}}$ are Z3-formalizable pre/post-condition assertions; and $\mathcal{S}$ documents applicable scope. Because $\phi_{\text{pre}}$ and $\phi_{\text{post}}$ are SMT-LIB expressions, each paradigm carries a mechanically checkable interface absent from prior natural-language memory formats. Screening is sample-based (Section~\ref{ssec:verification}), providing probabilistic—not formal—correctness certificates.

\textbf{Position in the inference-rule spectrum.} PRISM paradigms occupy the gap between (i) EBL macro-operators~\citep{MitchellKellerKedarCabelli1986}, which carry provable soundness but require a hand-built domain theory, and (ii) unverified LLM-memory heuristics, which need no domain theory but offer no mechanical quality control. Paradigms are cross-instance reusable (unlike single-invocation SMT lemmas), LLM-induced from clustered traces (unlike hand-crafted EBL schemas), and empirically Z3-screened before storage (unlike prior LLM-memory units).

\subsection{PRISM: Architecture Overview}
\label{ssec:overview}

PRISM operates via two decoupled stages.

\textbf{Offline stage (run once per task distribution).} Given $N$ training puzzles with ground-truth solutions:
\begin{enumerate}
    \item \textbf{Trajectory collection:} Run LLM+Z3 on training puzzles; retain only successful trajectories.
    \item \textbf{KDP identification:} Extract Key Decision Points from each trajectory (Section~\ref{ssec:kdp}).
    \item \textbf{Clustering:} Group similar KDPs by constraint-type feature vectors via agglomerative clustering.
    \item \textbf{Paradigm abstraction:} Synthesize one paradigm per cluster via LLM.
    \item \textbf{Triple Z3 verification:} Screen candidates along three axes (Section~\ref{ssec:verification}).
    \item \textbf{Library ingestion:} Persist verified paradigms in Paradigm Library $\mathcal{L}$.
\end{enumerate}
Result: $\approx$35 verified paradigms extracted from $\approx$1,500 trajectories (40:1 compression ratio).

\textbf{Online stage (per puzzle).} For each new puzzle $\mathcal{P}^*$:
\begin{enumerate}
    \item \textbf{Initial translation:} LLM formalizes natural-language constraints to Z3 assertions.
    \item \textbf{Paradigm retrieval:} Extract constraint-type signature; retrieve via two-layer filter (Section~\ref{ssec:retrieval}).
    \item \textbf{Consistency pre-check and hint injection:} Verify paradigm operations against current solver state before injection.
    \item \textbf{Step verification:} Z3 verifies each LLM inference step.
    \item \textbf{Repair loop:} On UNSAT, extract UNSAT core, record typed repair in $\mathcal{M}$, detect stagnation/loops.
    \item \textbf{Strategy escalation:} Trigger four-level protocol (Section~\ref{ssec:escalation}).
    \item \textbf{Write-back:} Stage successful repairs for batch re-verification and promotion to $\mathcal{L}$ (Section~\ref{ssec:writeback}).
\end{enumerate}

\subsection{Offline Stage: Paradigm Distillation}

\subsubsection{Trajectory Collection}

We execute LLM+Z3 solving on $N = 600$ training puzzles with generation temperature $T = 0.7$, running $r = 3$ independent generations per puzzle. Only successful trajectories (Z3 returns SAT; assignment matches ground truth) are retained, yielding $\approx$1,500 successful trajectories. Each step $a_t$ records: Z3 verification outcome (SAT/UNSAT/ERROR), domain change record $\Delta_t$, constraint-type signature, and step-type classification.

\subsubsection{Key Decision Point Identification}
\label{ssec:kdp}

A step $(S_{t-1}, a_t, S_t)$ qualifies as a Key Decision Point if it satisfies at least one of four conditions:

\textbf{Condition A (Domain reduction):} $\exists x \in \mathbf{X}: |\delta_{t-1}(x)| - |\delta_t(x)| \geq 2$. At least one variable's domain shrinks by two or more in a single step—a decisive, non-trivial inference.

\textbf{Condition B (Non-trivial step type):} The step is classified as \textsc{Chain} (transitive propagation) or \textsc{Contradiction} (explicit elimination of infeasible assignments). Classification uses lightweight rule-based parsing of the natural-language action text.

\textbf{Condition C (Aggregate information gain):} The aggregate log-ratio of domain reductions across all variables exceeds a soft bit-threshold $\tau_{\text{ig}} = 1.0$:
\[
G(t) = \sum_{\substack{x:\; |\delta_t(x)| < |\delta_{t-1}(x)|}} \log_2 \frac{|\delta_{t-1}(x)|}{|\delta_t(x)|} \;\geq\; \tau_{\text{ig}}.
\]
Condition C complements the hard per-variable cutoff of A by capturing \textit{gradual tightening} steps—reducing several domains from $5\!\to\!3$ or $4\!\to\!2$—that are individually below the A threshold but collectively constitute a meaningful inference. This pattern dominates harder puzzles (e.g., 6$\times$6 ZebraLogic) and substantially broadens paradigm coverage for those scales.

\textbf{Condition D (Successful repair):} A repair step $a_t$ (action type = \textsc{repair}) that causes the solver to return SAT in a trajectory ultimately marked successful. This ensures that effective repair patterns—which differ structurally from forward inference—are captured as KDPs and enter the distillation pipeline.

Across $\approx$1,500 trajectories, these four conditions extract $\approx$6,000 KDPs.

\subsubsection{Cross-Trajectory Clustering and Paradigm Abstraction}

\textbf{Feature representation.} Each KDP is represented as a fixed-dimensional vector $\mathbf{f}(\text{kdp}) = [\mathbf{h}_{\text{ct}}, \mathbf{b}_{\text{dom}}, \mathbf{e}_\tau, n_{\text{var}}, n_{\text{con}}]$, where $\mathbf{h}_{\text{ct}}$ is a normalized histogram over 9 constraint types (direct/relative position, adjacency, exclusion, inclusion, ordering, binding, all-different, logical-implication); $\mathbf{b}_{\text{dom}}$ is a 4-bin domain-size distribution; $\mathbf{e}_\tau$ is a step-type one-hot encoding; and $n_{\text{var}}, n_{\text{con}}$ are normalized problem-scale metrics.

\textbf{Clustering.} We apply complete-linkage agglomerative clustering with cosine distance and threshold $\theta = 0.25$ (selected via grid search on a held-out development set). Clusters with fewer than $m_{\min} = 5$ members are discarded.

\textbf{Paradigm abstraction.} For each qualifying cluster, we uniformly sample up to 10 representative KDPs (solver state, reasoning text, domain-change indicators) and prompt an LLM to synthesize a structured paradigm: \textit{trigger conditions, operation template, Z3 pre/post-conditions, and scope}. Up to 2 retries per cluster accommodate LLM variability.

\subsubsection{Z3 Triple Verification}
\label{ssec:verification}

Each candidate paradigm $\hat{P}$ undergoes three independent checks; failure in any one results in rejection.

\textbf{(1) Soundness.} The paradigm's operation must be consistent with random partial views of its precondition. We sample $N_p = 50$ random puzzle states satisfying $\hat{P}$'s trigger conditions, instantiate the operation, and test $\text{Z3-SAT}(C_i \cup \{\text{op}(\hat{P})\})$ for each. Acceptance requires $\text{soundness}(\hat{P}) \geq \tau_{\text{snd}} = 0.90$.

\textbf{(2) Effect.} The operation must genuinely constrain the solution space—it must be neither contradictory nor vacuously implied by the precondition. Two sub-checks must both hold:
\begin{itemize}
    \item \textbf{Non-contradiction:} $\text{Z3-SAT}(\phi_{\text{pre}} \wedge \text{op}(\hat{P}))$ — the operation is consistent with the precondition.
    \item \textbf{Non-vacuousness:} $\text{Z3-SAT}(\phi_{\text{pre}} \wedge \neg\,\text{op}(\hat{P}))$ — there exists a state where the precondition holds but the operation does not. Equivalently, the precondition does not already entail the operation. If this check fails ($\phi_{\text{pre}} \wedge \neg\,\text{op}$ is UNSAT), the paradigm is tautological under its trigger and adds no information.
\end{itemize}

\textbf{(3) Trigger precision.} The trigger conditions must be selective. We sample $N_p = 50$ random subsets of size 2--5 from a representative constraint pool of 18 annotated entries spanning ZebraLogic-domain and generic integer-arithmetic patterns. Each entry is tagged with canonical constraint kinds from a fixed taxonomy $\mathcal{K}$ (10 kinds, including \texttt{position\_arithmetic}, \texttt{attribute\_binding}, \texttt{distinct}, etc.). A paradigm's trigger fires on sample $s$ iff $\mathcal{K}_{\hat{P}} \cap \mathcal{K}_s \neq \emptyset$ (set-intersection, not substring matching—preventing spurious lexical co-occurrences). Acceptance requires a minimum non-firing rate $\tau_p = 0.20$: the trigger must \textit{not} fire on at least 20\% of random samples, ensuring the paradigm is not pathologically over-general. Paradigms accepted on all three checks receive a confidence score (equal to their soundness rate) and are persisted in Paradigm Library $\mathcal{L}$.

\subsection{Online Stage: Paradigm-Guided Inference and Adaptive Repair}

\subsubsection{Two-Layer Paradigm Retrieval}
\label{ssec:retrieval}

At each step $t$, we extract the constraint-type signature $\sigma(S_t) = \text{sorted}(\{type(c) \mid c \in C_t\})$ and apply:

\textbf{Layer 1 (Type matching).} Retrieve paradigms $P \in \mathcal{L}$ whose declared constraint kinds $\mathcal{K}_P$ intersect with $\sigma(S_t)$, filtered by confidence floor $\tau_{\text{conf}} = 0.50$. To address multiplicity (two puzzles with the same \textit{set} of constraint types may differ in \textit{count}), each paradigm's trigger may declare optional relational predicates from a small vocabulary: \texttt{count\_atleast}($t, n$), \texttt{count\_atmost}($t, n$), \texttt{cooccur}($t_1, t_2$). Predicates are evaluated against the actual type bag (a multiset); paradigms failing their predicates are pruned before ranking. Layer 1 typically reduces candidates from $\approx$35 to $<$5 in $O(|\mathcal{L}|)$ time.

\textbf{Layer 2 (LLM semantic judgment, cost-aware).} When $|\mathcal{L}_1(S_t)| \geq 2$, an additional LLM call judges which retrieved paradigms have trigger conditions genuinely satisfied by the current state. Under the default \texttt{complexity\_gated} policy, Layer 2 activates only when (i) at least 2 candidates survive Layer 1 \textit{and} (ii) either the puzzle has $\geq$25 constraints or at least one UNSAT has been observed. This bounds Layer-2 cost to high-complexity or repair-phase steps.

\subsubsection{Consistency Pre-Check and Hint Injection}

Before injection, each surviving paradigm $P$ undergoes a Z3 consistency pre-check:
\[
\text{consistent}(P, S_t) = \text{Z3-SAT}(C_t \cup \{\text{op}(P)\}).
\]
Only paradigms passing this check are injected as soft (advisory) hints; they are not added as hard constraints. This second safety barrier filters paradigms whose operations conflict with constraints added since the paradigm was trained.

\subsubsection{Step Verification and UNSAT Attribution}

After the LLM generates inference $a_t$, Z3 verifies it:

\textbf{SAT:} Update domain assignments; save the resulting state as a potential revert checkpoint. Continue solving.

\textbf{UNSAT:} Extract UNSAT core $U_t$ (the minimal unsatisfiable constraint subset) and perform two-level causal attribution. \textit{Primary attribution} is coarse: if $a_t \in U_t$, the new assertion is responsible (\textsc{New\_Assertion}); if $a_t \notin U_t$, the contradiction predates $a_t$ (\textsc{Legacy\_Error}). \textit{Secondary classification} refines \textsc{New\_Assertion} into one of four structured \textit{ErrorTypes} via AST parsing rather than substring matching, making the classifier robust to Z3's equivalent surface forms: \textbf{SyntaxError}, \textbf{ScopeError} (undefined identifier), \textbf{SemanticFlip} (contradictory comparison operators on identical operands, detected via polarity-tracking AST walk), or \textbf{OverConstraint} (default: the new assertion legally tightens the formula beyond satisfiability). The ErrorType is stored in the repair record and used by the escalation logic to select the most appropriate repair class.

\subsection{Repair Trajectory Memory}
\label{ssec:memory}

\subsubsection{Data Structure}

Repair Trajectory Memory $\mathcal{M}$ maintains ordered records $\{R_1, \ldots, R_L\}$, each a six-tuple:
\[
R_t = (e_t,\; U_t,\; \hat{U}_t,\; \alpha_t,\; o_t,\; U'_t)
\]
where $e_t$ is the error-type label; $U_t$ is the UNSAT core (constraint IDs); $\hat{U}_t = \text{SHA-256}(\text{sorted}(\tilde{U}_t))$ is a canonicalized core fingerprint enabling rapid cross-iteration comparison; $\alpha_t$ is a structured repair action triple $(\tau, c, \pi)$ (type, target constraint, parameter signature) plus a sentence-transformer embedding (all-MiniLM-L6-v2, 384-dim); $o_t \in \{\text{SAT}, \text{UNSAT}\}$ is the repair outcome; and $U'_t$ is the new UNSAT core post-repair if $o_t = \text{UNSAT}$.

\textbf{Core canonicalization.} Z3's \texttt{unsat\_core()} is not unique—different invocations on the same logically unsatisfiable set may return different cores depending on assumption ordering. Naive comparison can therefore both under-detect stagnation (same underlying contradiction surfaces as different core sets) and over-detect it (two unrelated contradictions share a high-frequency constraint such as \texttt{Distinct(...)}). PRISM canonicalizes each core via (i) operator aliasing (e.g., $\texttt{Not}(x > 0) \leftrightarrow x \leq 0$) and (ii) variable renaming to first-occurrence order. The fingerprint $\hat{U}_t$ is the SHA-256 of the normalized, sorted canonical core.

\subsubsection{Stagnation Detection}

Stagnation is detected when the last $k = 3$ repair records show high UNSAT core overlap:
\[
\text{stagnate}(\mathcal{M}, k{=}3) = \mathbf{1}\!\left[\,\max_{i \neq j \in \text{last-}k}\; J(\tilde{U}_i, \tilde{U}_j) \geq \tau_{\text{stag}}\right]
\]
where Jaccard similarity $J(A, B) = |A \cap B| / |A \cup B|$ is computed over canonicalized cores and $\tau_{\text{stag}} = 0.75$. The \textit{maximum} (any-pair) rather than minimum (all-pairs) is used as the trigger criterion: early-warning detection preserves more of the repair budget for recovery.

\subsubsection{Loop Detection}

Loop detection identifies when a proposed repair duplicates a historical attempt, combining two complementary signals:

\textbf{(i) Structured signal (primary).} Each repair action carries a triple $(\tau, c, \pi)$: repair type (drawn from a closed vocabulary of 8 classes, e.g., \texttt{relax\_bound}, \texttt{split\_constraint}), target constraint ID, and parameter signature. A loop is flagged when this triple exactly matches any prior record, regardless of natural-language phrasing.

\textbf{(ii) Semantic signal (fallback).} When the structured triple is absent or yields no match, cosine similarity between sentence-transformer embeddings of action summaries is used:
\[
\text{loop}(\alpha_{\text{new}}) = \mathbf{1}\!\bigl[\bigl((\tau, c, \pi)_{\text{new}} = (\tau, c, \pi)_t \text{ for some }t\bigr)\;\vee\;\bigl(\max_t \cos(\mathbf{e}_{\text{new}}, \mathbf{e}_t) \geq \tau_\ell\bigr)\bigr]
\]
with $\tau_\ell = 0.90$. The structured-signal-first design detects loops even when the LLM paraphrases the same action (which routinely produces cosines in the $0.78$--$0.85$ range, below the safe semantic threshold).

\subsubsection{Four-Level Strategy Escalation}
\label{ssec:escalation}

Upon stagnation or loop detection, PRISM escalates through four levels in order of increasing disruption:

\begin{enumerate}
    \item \textbf{L1 — Switch target.} Preserve the current repair type but target a different constraint from the UNSAT core. Triggered when the same constraint has been targeted in consecutive repair rounds.
    \item \textbf{L2 — Switch type.} Change the repair strategy class (e.g., \texttt{relax\_bound} $\to$ \texttt{split\_constraint}). Triggered when the same repair type dominates the recent history window.
    \item \textbf{L3 — Revert checkpoint.} Pop the most recent SAT-adjacent state from a bounded LIFO stack (depth $\leq 5$) and roll back the solver. Successive L3 escalations pop progressively older checkpoints; a depleted stack falls through to L4. Triggered when stagnation is detected and a checkpoint is available.
    \item \textbf{L4 — Full retranslation.} Discard the current formalization; re-translate all constraints from the original natural-language description, using the accumulated repair history as guidance.
\end{enumerate}

\textbf{Termination guarantee.} The online phase terminates within a bounded LLM-call budget. The repair loop is capped at $R = 5$ rounds; L4 is triggered at most once per puzzle (a second L4 condition is suppressed, marking the puzzle unsolved); and the L3 checkpoint stack has depth $\leq 5$. The worst-case per-puzzle budget is $1\;(\text{translate}) + R\;(\text{repair rounds}) + 1\;(\text{L4 retranslate}) = R + 2 = 7$ LLM calls. A formal proof sketch is given in Appendix~B.

\subsection{Paradigm Write-Back and Library Maintenance}
\label{ssec:writeback}

When a repair or forward inference step succeeds, PRISM does not immediately admit the corresponding pattern to $\mathcal{L}$. A \textit{staged write-back} policy is used:

\begin{enumerate}
    \item \textbf{Candidate staging (online, $O(1)$).} The responsible (trigger features, operation, inferred pre/post-condition) tuple is appended to a candidate pool $\mathcal{L}^{\dagger}$. No verification is performed online, preserving low latency.
    \item \textbf{Batched re-verification (offline, every $K = 50$ solved puzzles).} Candidates in $\mathcal{L}^{\dagger}$ are de-duplicated against $\mathcal{L}$ (by trigger-kind signature and operation-template hash) and re-run through the same triple verification protocol (Section~\ref{ssec:verification}).
    \item \textbf{Promotion.} Candidates passing all three checks are admitted to $\mathcal{L}$; failures are discarded with provenance retained. When an online step matches an existing $P \in \mathcal{L}$ and succeeds, only its $\langle\text{support\_count}, \text{confidence}\rangle$ counters are updated atomically—no re-verification required.
\end{enumerate}

The candidate pool $\mathcal{L}^{\dagger}$ is not used for retrieval during online inference, so $\mathcal{L}$ always maintains the same SMT-screened quality guarantee.

\subsection{Computational Cost}
\label{ssec:cost}

PRISM's cost has two components that amortize differently.

\textbf{Offline (one-time, amortized over the test set).} LLM calls arise in three places: trajectory collection ($N_{\text{train}} \cdot r \cdot \bar{T}$, where $\bar{T}$ is mean trajectory length); paradigm abstraction ($\leq 2|\mathcal{C}|$ calls for $|\mathcal{C}|$ clusters with up to 2 retries); and batched write-back re-verification. Z3 work scales as $N_p \cdot |\mathcal{P}_{\text{candidate}}|$ verification trials. The amortized offline cost per test puzzle is $C_{\text{offline}} / N_{\text{test}}$, which becomes negligible for large test sets.

\textbf{Online (per puzzle, worst case).} Per the termination guarantee, the online phase invokes the LLM at most $R + 2 = 7$ times per puzzle: one translation, at most $R = 5$ repair-round inferences, and at most one L4 retranslation. Layer-2 semantic matching adds at most one additional LLM call per step (gated by the \texttt{complexity\_gated} policy, contributing $\leq R$ additional calls on complex instances, zero on easy ones). Z3 work per step is bounded to one consistency pre-check, one solver call, and one UNSAT-core extraction.
